\section{Scope and architectural reset}

SDF-MPNEO vNext is a new mathematical architecture rather than an incremental
extension of the legacy fixed-grid Maxwell--thermal pipeline. The legacy
implementation remains useful as a numerical oracle for regression studies, but
its discretization choices are not part of the vNext definition.

The target problem is a coupled electromagnetic--thermal multi-port system with
continuously varying geometry, arbitrary rigid poses, finite conductor
cross-sections, packages and nearby material bodies, and time-dependent
electrical excitation. The model must support rapid repeated inference over a
large geometric/material design space while retaining an independently
checkable path to high accuracy.

\begin{principle}[Continuous scene first]
The primary object is a continuous physical scene. A mesh, point cloud,
quadrature rule, panelization, or basis truncation is an implementation detail
of a teacher or correction backend; it is never part of the public geometry
semantics.
\end{principle}

\begin{principle}[Learn operators, not discretized fields]
Neural networks approximate low-dimensional maps between physical operators,
latent modal coefficients, and continuous query functions. They do not learn a
vector whose coordinates are tied to cells of one background grid.
\end{principle}

\begin{principle}[Hard structure before soft loss]
Passivity, reciprocity, positive dissipation, non-negative heat generation, and
stable thermal dynamics are imposed by parameterization whenever possible.
Penalty terms are reserved for properties that cannot be enforced exactly.
\end{principle}

\begin{principle}[Fast prediction and certified correction are separate modes]
A purely learned/analytic prediction path is optimized for latency. A second
path evaluates physical integral-operator residuals and performs a bounded
number of corrections. High accuracy therefore does not require the neural
model to be uniformly exact everywhere.
\end{principle}

The intended online complexity is controlled by the number of ports, compact
electromagnetic modes, material objects, thermal states, and requested spatial
query points. It must not scale with the volume of an artificial far-field box.
